\documentclass[12pt]{article}
\usepackage[margin=1in]{geometry}
\usepackage{titlesec}
\usepackage{hyperref}
\usepackage{booktabs}
\usepackage{longtable}
\usepackage{parskip}
\usepackage{enumitem}
\usepackage{fancyhdr}

\hypersetup{colorlinks=true, linkcolor=blue, urlcolor=blue}

\pagestyle{fancy}
\fancyhf{}
\rhead{Campus Resource Booking \& Check-In System}
\lhead{User Manual}
\rfoot{Page \thepage}

\title{
    \vspace{2cm}
    {\Large User Manual} \\[1em]
    {\LARGE \textbf{Campus Resource Booking \& Check-In System}} \\[2em]
    {\large CS3338 --- Group 1} \\[0.5em]
    {\normalsize Version 4.0}
    \vspace{2cm}
}
\author{Ifunanya Okafor}
\date{May 15, 2025}

\begin{document}

\maketitle
\thispagestyle{empty}
\newpage

\tableofcontents
\newpage

% -------------------------------------------------------
\section*{Versions Table}
\addcontentsline{toc}{section}{Versions Table}

\begin{longtable}{|p{1.4cm}|p{2.4cm}|p{3.5cm}|p{6.5cm}|}
\hline
\textbf{Ver.} & \textbf{Date} & \textbf{Author(s)} & \textbf{Changes} \\
\hline
1.0 & Feb 2, 2025  & Ifunanya Okafor & Initial draft. Project overview, why it matters, snapshot 1 objectives, GitHub setup instructions. \\
\hline
2.0 & Mar 2, 2025  & Ifunanya Okafor & Added login/register instructions, how to browse resources and create bookings, cancellation steps. Updated run instructions. \\
\hline
3.0 & Apr 6, 2025  & Ifunanya Okafor & Added check-in usage instructions, admin dashboard walkthrough, staff resource management guide. \\
\hline
4.0 & May 15, 2025 & Ifunanya Okafor & Final revision. Added complete formal objective breakdown, all 4 snapshot reflections, full prerequisite table, updated seed accounts. \\
\hline
\end{longtable}
\newpage

% -------------------------------------------------------
\section{Project Management}

\textbf{Jira Project Board:}
\url{https://YOUR-TEAM.atlassian.net/jira/software/projects/CRB/boards}

\textit{Replace the URL above with your team's actual Jira link before submission.}

\textbf{GitHub Repository:}
\url{https://github.com/nanYaIfu/CS3338-Group1-Final-Project}

% -------------------------------------------------------
\section{Why This Software Matters}

College campuses rely on shared physical resources --- study rooms, computer labs, tutoring appointments --- that are typically managed through informal systems: paper sign-up sheets, email chains, or first-come-first-served walk-ins. This creates predictable problems:

\begin{itemize}
    \item Students waste time walking to spaces that are already occupied.
    \item No-shows waste reserved slots that other students could have used.
    \item Staff have no way to track demand, identify peak hours, or plan capacity.
    \item There is no accountability: anyone can claim or abandon a space without record.
\end{itemize}

The Campus Resource Booking \& Check-In System addresses all of these by providing:
\begin{itemize}
    \item A centralized, role-aware web platform where students can see real-time availability.
    \item A booking and check-in system that creates accountability for reservations.
    \item An admin analytics dashboard that gives staff data to make resource allocation decisions.
\end{itemize}

% -------------------------------------------------------
\section{Formal Objective Breakdown}

\subsection{Snapshot 1 --- Feb 2, 2025 (Foundation)}
\textbf{Goal:} Establish the technical foundation and complete all documentation scaffolding.
\begin{itemize}
    \item Define the project scope, MVP features, and technology stack
    \item Set up GitHub repository with folder structure and branch conventions
    \item Create and configure Jira project with Snapshot 1 sprint
    \item Author initial versions of: SDD, SRS, Design Spec, User Manual, Snapshot Objectives
    \item Create \texttt{docker-compose.yml} with frontend, backend, and MongoDB services
    \item Create high-level workflow diagram
\end{itemize}

\subsection{Snapshot 2 --- Mar 2, 2025 (Core Features)}
\textbf{Goal:} Implement authentication and the full resource booking flow.
\begin{itemize}
    \item Build user registration and login with JWT-based authentication
    \item Implement role-based access control (student / staff / admin)
    \item Build Resource CRUD endpoints and the resource browse/filter page
    \item Build booking creation with server-side conflict detection
    \item Build My Bookings page with cancel functionality
    \item Run and document TestRail test suite (Snapshot 2 report)
    \item Update SDD and SRS with implementation details
\end{itemize}

\subsection{Snapshot 3 --- Apr 6, 2025 (Advanced Features)}
\textbf{Goal:} Add check-in system and admin dashboard.
\begin{itemize}
    \item Implement check-in endpoint with time-window validation (15 min before / 10 min after)
    \item Add no-show detection via background interval job
    \item Build admin analytics dashboard with usage stats and charts
    \item Build staff resource management UI (create/edit/deactivate)
    \item Run and document TestRail test suite (Snapshot 3 report)
    \item Update all documentation
\end{itemize}

\subsection{Snapshot 4 --- May 15, 2025 (Finalization)}
\textbf{Goal:} Polish, bug fixes, and final documentation.
\begin{itemize}
    \item Resolve all bugs identified in Snapshot 2 and 3 TestRail runs
    \item Add rate limiting to authentication routes
    \item Polish UI responsiveness and error handling across all pages
    \item Confirm clean end-to-end Docker deployment
    \item Run and document final TestRail test suite (Snapshot 4 report)
    \item Finalize all documentation with complete version history
    \item Document future work and known limitations
\end{itemize}

% -------------------------------------------------------
\section{Prerequisites}

\begin{longtable}{|p{4cm}|p{3cm}|p{7cm}|}
\hline
\textbf{Tool} & \textbf{Version} & \textbf{Download} \\
\hline
Node.js & 18.x or later & \url{https://nodejs.org} \\
\hline
npm & 9.x or later & Bundled with Node.js \\
\hline
Docker Desktop & Latest stable & \url{https://docs.docker.com/get-docker} \\
\hline
Git & Any recent & \url{https://git-scm.com} \\
\hline
MongoDB (optional) & 6.x & Only needed if running without Docker \\
\hline
\end{longtable}

% -------------------------------------------------------
\section{How to Run the Project}

\subsection{With Docker (Recommended)}
\begin{enumerate}
    \item Clone the repository:
\begin{verbatim}
git clone https://github.com/nanYaIfu/CS3338-Group1-Final-Project.git
cd CS3338-Group1-Final-Project/campus-resource-booking-checkin
\end{verbatim}

    \item Copy and configure the environment file:
\begin{verbatim}
cp backend/.env.example backend/.env
# Open backend/.env and set a strong JWT_SECRET value
\end{verbatim}

    \item Start all services:
\begin{verbatim}
docker-compose up --build
\end{verbatim}

    \item Access the app:
\begin{verbatim}
Frontend:  http://localhost:3000
API:       http://localhost:5000
\end{verbatim}

    \item (Optional) Seed the database with test data:
\begin{verbatim}
docker exec crbs_backend node seed/seedData.js
\end{verbatim}
\end{enumerate}

\subsection{Without Docker}

\textbf{Backend:}
\begin{verbatim}
cd backend
npm install
npm run dev        # uses nodemon for hot-reload
\end{verbatim}

\textbf{Frontend (in a separate terminal):}
\begin{verbatim}
cd frontend
npm install
npm start
\end{verbatim}

Ensure MongoDB is running locally on \texttt{localhost:27017}, or set \texttt{MONGO\_URI} in \texttt{backend/.env} to a MongoDB Atlas connection string.

% -------------------------------------------------------
\section{Using the Application}

\subsection{Registering and Logging In}
\begin{enumerate}
    \item Navigate to \texttt{http://localhost:3000}.
    \item Click \textbf{Register} and fill in your name, email, password, and role.
    \item After registration you will be redirected to the Login page.
    \item Enter your credentials and click \textbf{Login}.
\end{enumerate}

\subsection{Browsing and Booking a Resource}
\begin{enumerate}
    \item After login, you will see the \textbf{Resource Browse} page.
    \item Use the type and location filters to narrow results.
    \item Click \textbf{Book} on any resource card.
    \item Select a date and an available time slot, then click \textbf{Confirm Booking}.
\end{enumerate}

\subsection{Checking In}
\begin{enumerate}
    \item Navigate to \textbf{My Bookings}.
    \item Find the upcoming booking you want to check in to.
    \item The \textbf{Check In} button will become active within 15 minutes of the booking start time.
    \item Click \textbf{Check In} to confirm your arrival. The status will update to ``Checked In.''
\end{enumerate}

\subsection{Cancelling a Booking}
\begin{enumerate}
    \item Navigate to \textbf{My Bookings}.
    \item Find the booking you wish to cancel.
    \item Click \textbf{Cancel}. You will be asked to confirm.
\end{enumerate}

% -------------------------------------------------------
\section{Default Seed Accounts}

\begin{longtable}{|p{2.5cm}|p{4cm}|p{3cm}|p{4cm}|}
\hline
\textbf{Role} & \textbf{Email} & \textbf{Password} & \textbf{Notes} \\
\hline
Admin   & admin@campus.edu   & admin123   & Development seed only \\
\hline
Staff   & staff@campus.edu   & staff123   & Development seed only \\
\hline
Student & student@campus.edu & student123 & Development seed only \\
\hline
\end{longtable}

\textit{Change all seed passwords before any public or production deployment.}

\end{document}

